\section{Structures de Conteneurs (Containers)}

Les conteneurs permettent d'organiser les widgets dans l'interface. L'API propose des structures pour \texttt{GtkBox}, \texttt{GtkGrid} et \texttt{GtkStack}.

\subsection{Structure \texttt{box\_config}}

\begin{lstlisting}[language=C]
typedef struct {
    GtkOrientation orientation;
    int spacing;
    bool homogeneous;
    widget_style_config style;
} box_config;
\end{lstlisting}

\begin{center}
\begin{tabular}{|l|l|p{5cm}|}
\hline
\textbf{Champ} & \textbf{Type} & \textbf{Description} \\
\hline
orientation & GtkOrientation & \texttt{GTK\_ORIENTATION\_HORIZONTAL} ou \texttt{VERTICAL} \\
\hline
spacing & int & Espace en pixels entre les enfants \\
\hline
homogeneous & bool & Si \texttt{true}, tous les enfants ont la même taille \\
\hline
style & widget\_style\_config & Configuration de base du widget \\
\hline
\end{tabular}
\end{center}

\bigskip

Le champ \texttt{orientation} détermine le sens de lecture visuelle de l'interface. En mode \texttt{VERTICAL}, les widgets s'empilent comme les éléments d'un menu ou d'un formulaire. En mode \texttt{HORIZONTAL}, ils se rangent côte à côte, ce qui est idéal pour créer des barres d'outils ou des rangées de boutons de navigation.

Le champ \texttt{spacing} permet de définir un espacement fixe et uniforme entre chaque élément enfant du conteneur. Cela garantit une cohérence visuelle immédiate sans avoir à configurer manuellement les marges individuelles de chaque widget, rendant l'interface plus aérée et lisible.

Le champ \texttt{homogeneous} est un outil de design puissant : lorsqu'il est activé, il force tous les widgets enfants à occuper exactement la même largeur (ou hauteur). Cela permet de créer des grilles de boutons parfaitement symétriques, quel que soit le texte qu'ils contiennent, assurant un rendu professionnel et équilibré.

Le champ \texttt{style} permet d'appliquer la configuration commune (\texttt{widget\_style\_config}) à l'ensemble du conteneur, facilitant ainsi la gestion des marges globales ou l'application de classes CSS thématiques à tout un groupe de widgets.

\subsection{Structure \texttt{grid\_config}}

\begin{lstlisting}[language=C]
typedef struct {
    int column_spacing;
    int row_spacing;
    bool homogeneous_columns;
    bool homogeneous_rows;
    widget_style_config style;
} grid_config;
\end{lstlisting}

\begin{center}
\begin{tabular}{|l|l|p{5cm}|}
\hline
\textbf{Champ} & \textbf{Type} & \textbf{Description} \\
\hline
column/row\_spacing & int & Espace entre les colonnes et les lignes \\
\hline
homogeneous\_* & bool & Force les colonnes/lignes à être de taille égale \\
\hline
\end{tabular}
\end{center}

\bigskip

Les champs \texttt{column\_spacing} et \texttt{row\_spacing} permettent de créer une grille aérée. Contrairement à une Box, la grille gère l'espace de manière bidimensionnelle, ce qui est crucial pour aligner parfaitement des étiquettes à gauche avec des champs de saisie à droite, comme dans un panneau de configuration.

Les champs \texttt{homogeneous\_columns} et \texttt{homogeneous\_rows} garantissent que toutes les cellules de la grille ont la même taille. C'est particulièrement utile pour créer un clavier numérique ou une calculatrice où chaque touche doit avoir des dimensions identiques pour conserver une esthétique régulière.

\subsection{Structure \texttt{stack\_config}}
\begin{lstlisting}[language=C]
typedef struct {
    GtkStackTransitionType transition;
    guint duration_ms;
    bool vhomogeneous;
    bool hhomogeneous;
    widget_style_config style;
} stack_config;
\end{lstlisting}
\begin{center}
\begin{tabular}{|l|l|>{\raggedright\arraybackslash}p{6cm}|}
\hline
\textbf{Champ} & \textbf{Type} & \textbf{Description} \\
\hline
transition & GtkStackTransitionType & Type d'animation de transition \\
\hline
duration\_ms & guint & Durée de l'animation en millisecondes \\
\hline
vhomogeneous / hhomogeneous & bool & Force les pages à avoir la même hauteur/largeur \\
\hline
\end{tabular}
\end{center}

\bigskip

Le champ \texttt{transition} définit l'effet visuel lors du changement de page (glissement latéral, fondu enchaîné, etc.). Cela transforme un changement de contenu brutal en une navigation fluide et moderne, typique des applications mobiles.

Le champ \texttt{duration\_ms} contrôle la vitesse de cette animation. Une durée courte (environ 300ms) donne une impression de réactivité, tandis qu'une durée plus longue permet à l'utilisateur de bien comprendre la transition entre deux sections de l'application.

Les champs \texttt{vhomogeneous} et \texttt{hhomogeneous} permettent de décider si le conteneur doit toujours garder la taille de sa plus grande page. Si désactivé, la fenêtre se redimensionnera automatiquement pour s'adapter au contenu de la page actuelle, ce qui offre plus de flexibilité.


\subsection{Exemple d'utilisation (Box)}

\begin{lstlisting}[language=C]
#include "container.h"

void create_layout(GtkWidget *parent) {
    box_config my_box = {
        .orientation = GTK_ORIENTATION_VERTICAL,
        .spacing = 10,
        .homogeneous = false,
        .style = { .margin_top = 15, .margin_bottom = 15 }
    };
    GtkWidget *box = create_box(&my_box);
    // gtk_box_append(GTK_BOX(parent), box);
}
\end{lstlisting}